\chapter{系统架构设计}

\section{总体架构}

本文档采用 Mermaid 语法描述系统架构组件图：

\begin{verbatim}
```mermaid
flowchart TB
    %% Frontend Package
    subgraph Frontend["Frontend Package"]
        UI["[Vite/React UI]"]
    end

    %% Backend Package
    subgraph Backend["Backend Package"]
        API["[Express API Server]"]
    end

    %% Storage Package
    subgraph Storage["Storage Package"]
        DB["[Local JSON DB]"]
    end

    %% Interfaces
    UI -- "REST (HTTP/JSON)" --> API
    API -- "File I/O" --> DB
```
\end{verbatim}

系统整体采用前后端分离架构：前端基于 Vite 和 React 构建，负责提供响应式的用户界面与富交互体验；后端采用 Node.js 平台下的 Express.js 框架，提供标准化的 RESTful API 接口服务，处理业务逻辑与核心模块功能；数据存储层采用轻量级的本地 JSON 文件系统，实现数据的快速读写与持久化，以满足轻量级和免数据库安装的部署需求。各层之间职责清晰，通过标准协议（HTTP/REST, File I/O）进行松耦合的数据流转。

\section{技术选型}

\begin{table}[H]
\centering
\caption{技术选型表}
\begin{tabular}{|l|l|l|}
\hline
\textbf{层级} & \textbf{技术} & \textbf{版本} \\
\hline
前端框架 & React + TypeScript & 18.x \\
\hline
构建工具 & Vite & 5.x \\
\hline
后端框架 & Express.js & 4.x \\
\hline
PDF 渲染 & PDF.js & 3.x \\
\hline
TTS & Web Speech API / Azure Speech & — \\
\hline
\end{tabular}
\end{table}

\section{模块划分}

% TODO: 模块列表与职责
